\hypertarget{appendix-a9f83ad40792}{}\label{referencestart:appendix-a9f83ad40792}
\section{Color encoding and linear-light operations}\label{appendix-a9f83ad40792-color-encoding-and-linear-light-operations}

\textbf{Curriculum depth:} Developed · \textbf{First seen:} XR01 · \textbf{Latest development:} University courses

\subsection{Definition and mechanism}\label{appendix-a9f83ad40792-definition-and-mechanism}

Color primaries/white point and nonlinear transfer encoding are different. Physical intensity filtering uses the appropriate linear quantity; black/white averaging then sRGB encoding gives about 0.735, not encoded 0.5. Tone/gamut mapping and absolute display luminance need separate contracts.

\subsection{Prerequisites}\label{appendix-a9f83ad40792-prerequisites}

\begin{itemize}
\tightlist
\item
  \hyperref[appendix-215c2570199b]{Light spectra and radiometry} --- The light quantity and units must be established before its interaction or visual effect is modeled.
\item
  \hyperref[appendix-f414bd4ed05c]{Interpretation contract} --- A numeric sample requires a domain, units and encoding convention.
\end{itemize}

\subsection{Used by}\label{appendix-a9f83ad40792-used-by}

\begin{itemize}
\tightlist
\item
  \hyperref[appendix-b2c1846508bf]{Perceptual metrics and display calibration}
\item
  \hyperref[appendix-92d45cdf2afc]{Exposure and sensor response}
\end{itemize}

\subsection{Recall and next depth}\label{appendix-a9f83ad40792-recall-and-next-depth}

State the metric's encoding and peak before comparing images. Pending: standardized HDR and measured display calibration. Written coverage is not a claim of learner mastery or accepted numbered-curriculum progress.

\subsection{Source and navigation}\label{appendix-a9f83ad40792-source-and-navigation}

\href{../07-Vision-XR-and-Playable-Worlds/main.pdf\#nameddest=04-visual-computing-ii-ch042-encoded-rgb-must-not-masquerade-as-linear-light}{Expanded teaching explanation (Vision, XR, and Playable Worlds, Vision, XR, and Playable Worlds, Chapter 9)} · \hyperref[appendix-fd090d982a15-xr01---independent-seeing-and-representation-source]{Independent coverage} · \hyperref[appendix-4ad47bdd269a]{Dependency Map}

\subsection{University course development}\label{appendix-a9f83ad40792-university-course-development}

Extend the existing color note with why image-processing arithmetic must distinguish encoded values from light. A display encoding such as sRGB uses a nonlinear transfer curve so code values allocate precision in a way that suits perception and storage. Averaging or adding those encoded numbers does not generally average or add emitted light. Decode to linear-light RGB before blending, filtering, estimating radiance, or compositing, then encode for display. If two equal-area samples contain linear intensities \(0\) and \(1\), their physical average is \(0.5\); averaging their encoded sRGB values and treating that result as light yields a different, brighter answer. Color primaries and a white point also define the coordinate system of RGB, so a numeric triplet is incomplete without its color space. Keep scene-referred radiance, display-referred color, and perceptual appearance separate.

\href{../03-Geometry-Cameras-and-Image-Formation/main.pdf\#nameddest=03-visual-computing-i-ch004-color-perception-encoding-and-display}{Color perception, encoding, and display (Geometry, Cameras, and Image Formation, Geometry, Cameras, and Image Formation, Chapter 11)} · \href{../03-Geometry-Cameras-and-Image-Formation/main.pdf\#nameddest=03-visual-computing-i-ch004-light-measurement-and-radiometry}{Light measurement and radiometry (Geometry, Cameras, and Image Formation, Geometry, Cameras, and Image Formation, Chapter 11)} · \href{../03-Geometry-Cameras-and-Image-Formation/main.pdf\#nameddest=03-visual-computing-i-ch004-image-sampling-filtering-and-efficient-reconstruction}{Image sampling, filtering, and efficient reconstruction (Geometry, Cameras, and Image Formation, Geometry, Cameras, and Image Formation, Chapter 11)}
